\begin{abstract}
Large language models, vision-language models, and vision-language-action
models are now routinely placed inside robots that talk to, perceive, and act
around people, and inside the research pipelines that study those robots. The
human-robot interaction (HRI) community has no shared standard for evaluating
them. Practices imported from the wider AI community, such as model-as-judge
scoring or human validation of outputs, are applied unevenly, and they answer
different questions depending on what the model is doing in the robot: text
that a robot speaks directly to a person must be judged for relevance,
correctness, safety, and transparency, whereas an embedding that feeds a
downstream decision is evaluated much like a classifier. This article argues
that a single standard is neither achievable nor desirable, and that the unit
of evaluation should be the \emph{role} a model plays, not the model. We
distinguish nine roles, six inside the robot's interaction loop (perceiver,
representer, decider, speaker, behaviour author, actor) and three inside the
research pipeline (annotator, judge, simulator), and four levels at which
evidence can be gathered (component, system, interaction, deployment). For
each role we identify the evaluation tradition it inherits: human evaluation
and meta-evaluation from natural language processing, inter-rater reliability
and validity theory from psychology and content analysis, and trial-based
policy evaluation and safety analysis from robotics. From these we derive 24
best practices, ten that apply to every study and fourteen that are specific
to a role, scaled by a three-tier scheme that ties the required evidence to
who is exposed to the model's output and what can go wrong. We then describe
practices that are new or newly feasible, including ceiling-referenced
agreement, equivalence testing for substitution claims, generalisability
designs that treat model runs and prompts as raters, frozen-interaction
replay, and reference-corpus alignment, and we provide worked examples,
precision-planning tables for validation subsets and physical trials, and a
reporting checklist for authors and reviewers. The aim is a common vocabulary
that lets HRI studies using generative models be compared, replicated, and
trusted in proportion to the evidence they report.
\end{abstract}
